% Part: normal-modal-logic
% Chapter: sequent-calculus
% Section: rules-for-K

\documentclass[../../../include/open-logic-section]{subfiles}

\begin{document}

\olfileid[fr]{nml}{seq}{rul}

\olsection{Règles pour \Log{K}}

Les règles des connecteurs propositionnels usuels sont les mêmes que
dans le calcul des séquents usuel~$\Log{LK}$. Les axiomes sont également
les mêmes : tout séquent de la forme~$!A \Sequent !A$ est un axiome.

Pour \iftag{notprvBox,notprvDiamond}{l’opérateur modal}{les opérateurs modaux}
\iftag{prvBox}{\iftag{prvDiamond}{$\Box$ et }{$\Box$}}{}%
\iftag{prvDiamond}{$\Diamond$}{}, nous ajoutons
\iftag{notprvBox,notprvDiamond}{la règle suivante}{les règles suivantes} :
\[
  \iftag{prvBox}{\iftag{prvDiamond}{% <> and [] primitive
    \Axiom$\Gamma \fCenter \Delta, !A$
    \RightLabel{$\Box$}
    \UnaryInf$\Box\Gamma \fCenter \Diamond\Delta, \Box!A$
    \DisplayProof  
    \qquad
    \Axiom$!A, \Gamma \fCenter \Delta$
    \RightLabel{$\Diamond$}
    \UnaryInf$\Diamond!A, \Box\Gamma \fCenter \Diamond\Delta$
    \DisplayProof
}{% only Box primitive
\Axiom$\Gamma \fCenter !A$
\RightLabel{$\Box$}
\UnaryInf$\Box\Gamma \fCenter \Box!A$
\DisplayProof
}}{% only <> primitive
  \Axiom$!A \fCenter \Delta$
  \RightLabel{$\Diamond$}
  \UnaryInf$\Diamond!A \fCenter \Diamond\Delta$
  \DisplayProof
}
\]
Ici, \iftag{prvBox}{$\Box\Gamma$ désigne la suite obtenue à partir
de~$\Gamma$ en plaçant~$\Box$ devant chaque !!{formula}%
\iftag{prvDiamond}{ ; }{}}{}%
\iftag{prvDiamond}{$\Diamond\Delta$ désigne la suite obtenue à partir
de~$\Delta$ en plaçant~$\Diamond$ devant chaque !!{formula}}{}.
\iftag{notprvDiamond}{Dans la règle~$\Box$, le membre droit de la
prémisse doit contenir au plus une !!{formula}~$!A$. }{}%
\iftag{notprvBox}{Dans la règle~$\Diamond$, le membre gauche de la
prémisse doit contenir au plus une !!{formula}~$!A$. }{}%
Il est permis que \iftag{prvBox}{$\Gamma$\iftag{prvDiamond}{ ou }{}}{}%
\iftag{prvDiamond}{$\Delta$}{} soit vide ; dans ce cas, la partie
correspondante \iftag{prvBox}{$\Box\Gamma$\iftag{prvDiamond}{ ou }{}}{}%
\iftag{prvDiamond}{$\Diamond\Delta$}{} du séquent de conclusion est vide
également.

La restriction à une seule !!{formula}~$!A$ pour
\iftag{prvBox}{l'introduction de~$\Box$ à droite
\iftag{prvDiamond}{ ou }{}}{}%
\iftag{prvDiamond}{celle de~$\Diamond$ à gauche}{}
est nécessaire. Si l'on pouvait
\iftag{prvBox}{préfixer de~$\Box$ un nombre quelconque de
!!{formula}s à droite\iftag{prvDiamond}{ ou }{}}{}%
\iftag{prvDiamond}{préfixer de~$\Diamond$ un nombre quelconque de
!!{formula}s à gauche}{}, on pourrait !!{derive} :
\[
    \iftag{prvBox}{
      \Axiom$!A \fCenter !A$
      \RightLabel{\RightR{\lnot}}
      \UnaryInf$\fCenter !A, \lnot !A$
      \RightLabel{$\Box*$}
      \UnaryInf$\fCenter \Box!A, \Box\lnot !A$
      \doubleLine
      \RightLabel{\RightR{\lor}}
      \UnaryInf$\fCenter \Box!A \lor \Box \lnot !A$
      \DisplayProof}{}
    \iftag{notprvBox,notprvDiamond}{}{\qquad}
    \iftag{prvDiamond}{
      \Axiom$!A \fCenter !A$
      \RightLabel{\LeftR{\lnot}}
      \UnaryInf$\lnot !A, !A \fCenter$
      \RightLabel{$\Diamond*$}
      \UnaryInf$\Diamond\lnot !A,\Diamond!A \fCenter $
      \RightLabel{\RightR{\lnot}}
      \UnaryInf$\Diamond!A \fCenter \lnot\Diamond\lnot !A$
      \RightLabel{\RightR{\lif}}
      \UnaryInf$ \fCenter \Diamond!A \lif \lnot\Diamond\lnot !A$
      \DisplayProof}{}
  \]
Or, aucune des formules obtenues n'est valide dans~$\Log{K}$ :
\iftag{prvBox}{$\Box!A \lor \Box \lnot !A$\iftag{prvDiamond}{ et
$\Diamond!A \lif \lnot\Diamond\lnot !A$}{}}{$\Diamond!A \lif
\lnot\Diamond\lnot !A$}.

Si l'on autorisait, dans la prémisse, des !!{formula}s auxiliaires en
plus de~$!A$, tout en laissant
\iftag{prvBox}{la règle~$\Box$ préfixer de~$\Box$ seulement~$!A$ à droite
\iftag{prvDiamond}{, ou }{}}{}%
\iftag{prvDiamond}{la règle~$\Diamond$ préfixer de~$\Diamond$
seulement~$!A$ à gauche}{} (sans modifier les !!{formula}s auxiliaires),
on pourrait !!{derive} :
\[
  \iftag{prvBox}{
    \Axiom$!A \fCenter !A$
    \RightLabel{\RightR{\lnot}}
    \UnaryInf$\fCenter !A, \lnot !A$
    \RightLabel{\RightR{\Exchange}}
    \UnaryInf$\fCenter \lnot !A, !A$
    \RightLabel{$\Box*$}
    \UnaryInf$\fCenter \lnot!A, \Box !A$
    \doubleLine
    \RightLabel{\RightR{\lor}}
    \UnaryInf$\fCenter \lnot!A \lor \Box !A$
    \DisplayProof}{}
  \iftag{notprvBox,notprvDiamond}{}{\qquad}
  \iftag{prvDiamond}{
    \Axiom$!A \fCenter !A$
    \RightLabel{\LeftR{\lnot}}
    \UnaryInf$\lnot !A, !A \fCenter$
    \RightLabel{$\Diamond*$}
    \UnaryInf$\Diamond\lnot !A, !A \fCenter $
    \RightLabel{\RightR{\lnot}}
    \UnaryInf$!A \fCenter \lnot\Diamond\lnot !A$
    \RightLabel{\RightR{\lif}}
    \UnaryInf$ \fCenter !A \lif \lnot\Diamond\lnot !A$
    \DisplayProof}{}
\]
Or, aucune des formules obtenues n'est valide dans~$\Log{K}$ :
\iftag{prvBox}{$\lnot!A \lor \Box !A$ (équivalente à $!A
\lif \Box!A$)\iftag{prvDiamond}{ et $!A \lif \lnot\Diamond\lnot !A$}{}}{$!A \lif \lnot\Diamond\lnot !A$}.

\end{document}
